\documentclass[10pt]{article}
\usepackage[T1]{fontenc}
\usepackage[utf8]{inputenc}
\usepackage[french]{babel}
\usepackage{amsmath,amssymb,mathrsfs,mathpazo}
\usepackage{geometry}
\geometry{a4paper,margin=2cm}
\usepackage{graphicx,array,tabularx,tabularray,colortbl,multirow,multicol,booktabs,enumitem,pifont,ulem,fancybox,diagbox,cancel}
\usepackage{pstricks,pstricks-add,pst-plot,pst-tree,pst-node,pst-text,pst-func,pst-math,pst-bspline,pst-3dplot,pst-eucl,pst-solides3d}
\let\pscalculate\undefined
\usepackage{pst-bezier}
\usepackage{tikz,pgfplots,tkz-tab}
\usetikzlibrary{arrows,arrows.meta,patterns,calc,positioning,trees,decorations.markings,shapes.geometric,shapes.misc,shapes.arrows,intersections,decorations.pathmorphing,decorations.pathreplacing,backgrounds,shadows,babel}
\pgfplotsset{compat=1.18}
\usepgfplotslibrary{fillbetween}
\usepackage[np]{numprint}
\usepackage{listings,makecell,pas-tableur,twoopt}
\let\barre\undefined
\usepackage[pstricks]{tablvar}
\definecolor{darkgreen}{RGB}{0,100,0}
\providecommand{\red}{\color{red}}
\providecommand{\blue}{\color{blue}}
\providecommand{\green}{\color{green}}
\providecommand{\black}{\color{black}}
\providecommand{\white}{\color{white}}
\providecommand{\gray}{\color{gray}}
\providecommand{\cyan}{\color{cyan}}
\providecommand{\magenta}{\color{magenta}}
\providecommand{\yellow}{\color{yellow}}
\providecommand{\violet}{\color{violet}}
\providecommand{\orange}{\color{orange}}
\providecommand{\euro}{\texteuro}
\providecommand{\EUR}{\texteuro}
\providecommand{\ieme}{\textsuperscript{e}}
\providecommand{\degres}{\ensuremath{^\circ}}
\providecommand{\degree}{\ensuremath{^\circ}}
\providecommand{\vv}[1]{\overrightarrow{#1}}
\providecommand{\Coord}[1]{#1}
\providecommand{\decosix}{}
\providecommand{\starredbullet}{\ensuremath{\bullet}}
\pagestyle{empty}
\setlength{\parindent}{0pt}
\usepackage{adjustbox}\DeclareUnicodeCharacter{2236}{:}\DeclareUnicodeCharacter{2190}{\ensuremath{\leftarrow}}
\providecommand{\R}{\mathbb{R}}
\providecommand{\N}{\mathbb{N}}
\providecommand{\Z}{\mathbb{Z}}
\providecommand{\Q}{\mathbb{Q}}
\providecommand{\D}{\mathbb{D}}
\providecommand{\C}{\mathbb{C}}
\providecommand{\e}{{\mathrm{e}}}
\providecommand{\vect}[1]{\overrightarrow{#1}}
\providecommand{\vectt}[1]{\overrightarrow{\mathrm{#1}}}
\providecommand{\vvt}[1]{\overrightarrow{\mathrm{#1}}}
\providecommand{\np}[2][]{{#2~#1}}
\providecommand{\vv}[1]{{\overrightarrow{#1}}}
\renewcommand{\d}{{\mathrm{d}}}
\providecommand{\degres}{{^\circ}}
\providecommand{\degre}{{^\circ}}
\providecommand{\blue}{\color{blue}}
\providecommand{\red}{\color{red}}
\providecommand{\enskip}{\quad}
\providecommand{\llbracket}{\lbrack\!\lbrack}
\providecommand{\rrbracket}{\rbrack\!\rbrack}
\providecommand{\euro}{{\text{€}}}
\providecommand{\no}{{\text{n°}}}
\providecommand{\interval}[2]{{[#1;#2]}}
\providecommand{\intervFO}[2]{{[#1;#2[}}
\providecommand{\dotfill}{\ldots}

\providecommand{\ds}{\displaystyle}
\providecommand{\textding}[1]{#1}
\providecommand{\Oij}{\(\mathrm{O};\vec{\imath},\vec{\jmath}\)}
\providecommand{\Oijk}{\(\mathrm{O};\vec{\imath},\vec{\jmath},\vec{k}\)}
\providecommand{\Ouv}{\(\mathrm{O};\vec{u},\vec{v}\)}
\providecommand{\intervFF}[2]{\ensuremath{[#1;#2]}}
\renewcommand{\euro}{\ifmmode\text{\texteuro}\else\texteuro\fi}
\DeclareUnicodeCharacter{00B0}{\ifmmode^{\circ}\else\textdegree\fi}
\DeclareUnicodeCharacter{20AC}{\ifmmode\text{\texteuro}\else\texteuro\fi}
% Énoncés : Éducation nationale. Transcription et corrigés : APMEP, Denis Vergès et contributeurs. Découpage et classement : CoopMaths, Sébastien Lozano, Éric Elter et contributeurs. Adaptation de la présentation : Maths Explorer.
% CC BY-SA : https://coopmaths.fr/www/about/
% https://coopmaths.fr/alea/static/bac/2026/tex/bac_2026_06_sujet2_etrangers_3.tex
\begin{document}
\noindent\textit{APMEP / CoopMaths. CC BY-SA. Présentation adaptée par Maths Explorer.}\par\medskip
\subsection*{Partie A}
On note $f$ la fonction définie sur l'intervalle $\left[2 ;+\infty\right[$ par \[f(x)=\ln \left(3 x^{2}+2 x\right).\]

On admet que la fonction $f$ est dérivable sur l'intervalle $\left[2 ; +\infty \right[$ et on note $f'$ sa fonction dérivée.

\begin{enumerate}
	\item Démontrer que la fonction $f$ est strictement croissante sur l'intervalle $[2 ;+\infty[$.

	\item On note $g$ la fonction définie sur l'intervalle $[2 ;+\infty[$ par :
	\[g(x)=f(x)-x\]

	On admet que la fonction $g$ est strictement décroissante sur l'intervalle $[2 ;+\infty[$

	et que $\lim\limits_{x \to +\infty} g(x)=-\infty$.

	\begin{enumerate}
		\item Démontrer que l'équation $g(x)=0$ admet une unique solution $\alpha$ sur l'intervalle $[2 ;+\infty[$.

		\item Donner la valeur de $\alpha$ arrondie au centième.

		\item En déduire le tableau de signes de la fonction $g$ sur l'intervalle $[2 ;+\infty[$.
	\end{enumerate}
\end{enumerate}

\subsection*{Partie B}
Dans cette partie, les réponses pourront s'appuyer sur les résultats de la \textbf{partie A}.

On définit une suite $(a_{n})$ par son premier terme $a_{0}>0$ et pour tout entier naturel $n$,
\[a_{n+1}=\ln \left(3 {a_{n}}^{2}+2 a_{n}\right)\]

On étudie le cas où \quad $2 \leqslant a_{0} \leqslant \alpha$, \quad où $\alpha$ est l'unique solution de l'équation $g(x)=0$.

\begin{enumerate}
	\item Démontrer que, pour tout entier naturel $n$, on a $2 \leqslant a_{n} \leqslant \alpha$.

	\item Démontrer que la suite $(a_{n})$ est croissante.

	\item Démontrer que la suite $(a_{n})$ converge.

	\item Démontrer que la limite de la suite $(a_{n})$ est $\alpha$.
\end{enumerate}

\subsection*{Partie C}
Dans cette partie, on prend $a_{0}=2$.

La suite $(a_{n})$ est ainsi définie par $a_{0}=2$ et pour tout entier naturel $n$,
\[a_{n+1}=\ln \left(3 {a_{n}}^{2}+2 a_{n}\right)\]

On note $(b_{n})$ la suite définie par $b_{0}=10$ et pour tout entier naturel $n$,
\[b_{n+1}=\ln \left(3 {b_{n}}^{2}+2 b_{n}\right)\]

On admet que la suite $\left(b_{n}\right)$ est strictement décroissante et qu'elle converge vers $\alpha$.

\begin{enumerate}
	\item Justifier que, pour tout entier naturel $n$, on a \quad $a_{n} \leqslant b_{n}$.

	\item On considère le script ci-dessous écrit en langage Python.

\begin{center}
	\begin{tblr}{colspec = {X[6cm,font=\ttfamily]},vlines, hline{1,Z}={solid},rowsep={0pt}}
	from math import *\\
	def algo(p):\\
	\quad a = 2\\
	\quad b = 10\\
	\quad n = 0\\
	\quad while b - a > 10**(-p):\\
	\quad \quad a = log(3*a**2 + 2*a)\\
	\quad \quad b = log(3*b**2 + 2*b)\\
	\quad \quad n = n + 1\\
	\quad return (n,a)\\
\end{tblr}
\end{center}

On rappelle qu'en langage Python :

	\begin{itemize}[label=\textbullet~]
		\item la commande \texttt{log(c)} renvoie la valeur de $\ln (c)$;

		\item la commande \texttt{a**2} renvoie la valeur de $a^{2}$.
	\end{itemize}

	\begin{enumerate}
		\item Donner les valeurs renvoyées par l'instruction \texttt{algo(2)}.

		\emph{On arrondira si besoin les valeurs au millième.}

		\item Interpréter les valeurs renvoyées dans le contexte de l'exercice.
	\end{enumerate}
\end{enumerate}

\end{document}
